\documentclass[11pt]{article}
\usepackage{../style/forage}
\usepackage{needspace}

% All tables and figures in figures/ are generated by the four notebooks (01-04).
% Re-run them if the data or assumptions change, then recompile.

% keep headings with the text that follows them
\let\oldsubheading\subheading
\renewcommand{\subheading}[1]{\Needspace{6\baselineskip}\oldsubheading{#1}}
\let\oldheading\heading
\renewcommand{\heading}[1]{\Needspace{8\baselineskip}\oldheading{#1}}

% figure / table with a grey caption underneath
\newcommand{\figurebox}[3][\linewidth]{%
  \begin{center}\includegraphics[width=#1]{#2}\par
  {\small\color{gray}#3}\end{center}}
\newcommand{\tablebox}[2]{%
  \begin{center}\small\input{#1}\par\smallskip
  {\color{gray}#2}\end{center}}
\newcommand{\widetablebox}[2]{%
  \begin{center}\resizebox{\linewidth}{!}{\input{#1}}\par\smallskip
  {\small\color{gray}#2}\end{center}}

\begin{document}

% ================= Title page =================
\begin{titlepage}
  \vspace*{4cm}
  {\color{navy}\fontsize{30}{36}\selectfont Portfolio and Operations\\[2mm]
  Review}\par\medskip
  {\Large Kestrel Bay Capital: Market Insight, Risk Rules,\\
  Allocation and Trade Reconciliation}\par
  {\color{rulegray}\rule{\linewidth}{0.6pt}}\par\bigskip
  \begin{tabular}{@{}ll}
    \textbf{Client:}      & Kestrel Bay Capital (fictional) \\
    \textbf{Prepared by:} & Thiziri Amezza, Sales and Trading Analyst \\
    \textbf{Scope:}       & One year of market data, 3,000 trades \\
    \textbf{Date:}        & \today \\
  \end{tabular}
  \vfill
  {\small\color{gray}Capstone project of the Bank of America Global Markets
  Sales and Trading job simulation portfolio. The client, the market and the
  trades are synthetic, generated by \texttt{src/generate\_data.py} with a
  fixed seed; every table and figure is produced by the notebooks
  \texttt{01}--\texttt{04} in this folder. Simulated results describe the
  model, not a forecast.}
\end{titlepage}

% ================= Executive summary =================
\heading{Executive Summary}
Kestrel Bay Capital had a strong year: the portfolio returned
\textbf{+16.3\%} against a 6\% target, and its worst fall from a peak was
only about 3\% against a 10\% limit. A calm year, however, says little about
the risks ahead. This review tests the portfolio against 5,000 simulated
years, searches every allocation the client's rules allow, and builds a tool
for the client's operational concern, trade reporting.

\begin{tcolorbox}[enhanced, colback=badgebg, frame hidden, arc=2mm,
    borderline west={1.2mm}{0pt}{navy}, left=5mm, right=4mm, top=3mm, bottom=3mm]
\textbf{\color{navy}Recommendations}
\begin{enumerate}[leftmargin=1.5em, itemsep=3pt]
  \item \textbf{Make small changes to the allocation:} Financials +5, Energy
  +5, Real Estate \textminus10 points. Today's mix is already close to the best
  the rules allow; this change adds about 0.1 point of return a year and cuts
  the loss from a +100 bp rate rise from 1.9\% to 1.4\%.
  \item \textbf{Adopt risk-off triggers if the 10\% drawdown limit is a hard
  limit.} They cut the chance of breaching it by about a third (from 5.8\% to
  4.2\% of simulated years) at a cost of about 1 point of median return.
  \item \textbf{Replace the manual trade check with the rule-based
  reconciliation engine.} It found all 270 breaks with no false alarms and
  cuts analyst time from about 1,000 to 45 hours a year.
  \item \textbf{Do not rely on one year of data to measure rate risk.} The
  one-year estimate pointed the wrong way; scenario analysis should be used
  instead.
\end{enumerate}
\end{tcolorbox}

\infobox{\faClipboardList}{Client at a glance}{
  \item \textbf{Type:} mid-sized multi-family office, \$250m under management.
  \item \textbf{Allocation:} Technology 35\%, Healthcare 15\%, Financials
  10\%, Renewable Energy 15\%, Real Estate 10\%, Cash 15\%.
  \item \textbf{Objectives:} at least 6\% a year; maximum drawdown under
  10\%; at least 15\% in sustainable assets; faster, reliable trade reporting.
  \item \textbf{Rules:} no sector above 40\%; cash 5--20\%; no leverage.
}

% ================= 1. Market insight =================
\newpage
\heading{1. The Year in Review}

\figurebox[0.92\linewidth]{figures/s1_fig_market.pdf}{Figure 1 --- Sector
performance, VIX and policy rate over 52 weeks. The rate rose from 3.00\% to
3.75\% and was then cut to 3.25\%.}

\tablebox{figures/s1_table_sectors.tex}{Table 1 --- Sector statistics over the
year. ``Move per +25 bp'' is estimated from the data (see the caution below).}

\subheading{What the year shows}
\begin{itemize}[leftmargin=1.5em, itemsep=2pt]
  \item \textbf{Concentration:} Technology is 35\% of the capital but
  \textbf{51\% of the portfolio's risk} (Figure 2).
  \item \textbf{Sustainability pays:} Renewable Energy had the best return per
  unit of volatility (1.65), so the client's sustainability objective did not
  cost return.
  \item \textbf{A gap:} the client holds no Energy, which rose 25\% and tends
  to rise with interest rates --- a natural hedge for a portfolio that loses
  when rates go up.
  \item \textbf{A caution:} the policy rate moved in only 5 weeks, so the
  rate sensitivities in Table 1 rest on 5 observations. Section 2 shows they
  are badly wrong.
\end{itemize}

\figurebox[0.85\linewidth]{figures/s1_fig_risk.pdf}{Figure 2 --- Share of
capital vs.\ share of risk for each sector held.}

% ================= 2. Live dynamics =================
\heading{2. Testing the Risk Limits}
One calm year cannot show how often the 10\% drawdown limit would be
breached. I compared three ways of managing the portfolio:
\begin{itemize}[leftmargin=1.5em, itemsep=2pt]
  \item \textbf{Buy and hold:} set the allocation once and let it drift.
  \item \textbf{Monthly rebalance:} trade back to the target every 4 weeks.
  \item \textbf{Monthly rebalance + risk triggers:} also switch to a
  defensive allocation (less Technology, more Healthcare and cash, within the
  client's rules) when the VIX goes above 20 or the portfolio is 5\% below its
  peak; switch back when the VIX is below 17 and the drawdown under 3\%.
\end{itemize}
Each approach was replayed on the actual year and run through 5,000
simulated years drawn from the same market model, with random crises, rate
moves and sector shocks. Trading costs are included, and decisions use only
the previous week's data.

\tablebox{figures/s2_table_replay.tex}{Table 2 --- Replay of the actual year.}

\tablebox{figures/s2_table_mc.tex}{Table 3 --- Results across 5,000 simulated
years.}

\figurebox[0.9\linewidth]{figures/s2_fig_mc.pdf}{Figure 3 --- The trade-off:
the triggers reduce the chance of breaching the drawdown limit but raise the
chance of missing the return target.}

\subheading{What this means}
\begin{itemize}[leftmargin=1.5em, itemsep=2pt]
  \item The actual year cannot separate the approaches: all returned about
  +16\%.
  \item Over 5,000 years, the triggers cut the chance of a drawdown worse than
  10\% \textbf{from 5.8\% to 4.2\%}, but the median return falls from 13.3\%
  to 12.2\% and the chance of missing the 6\% target rises from 20.8\% to
  23.6\%. Which matters more is the client's decision; the simulation puts a
  number on it.
\end{itemize}

\subheading{Rate risk: the one-year estimate pointed the wrong way}
Because the market is synthetic, the true rate sensitivities are known.
Table 4 compares them with the estimates from the year's data.

\tablebox{figures/s2_table_rates.tex}{Table 4 --- Move per +25 bp: true model
vs.\ one-year estimate, and the portfolio impact of two rate shocks.}

With the true sensitivities, a +100 bp rise costs the portfolio
\textbf{1.9\%}; the one-year estimates predicted a small \emph{gain}. Rate
risk should be assessed with scenarios, not with a single year's regression.

% ================= 3. Allocation =================
\heading{3. Recommended Allocation and Products}

\subheading{Products}
Each building block is mapped to a product and the objective it serves.
Derivatives (for example a protective put on the technology ETF) were
considered but excluded: the client's rules rule out leverage, and the market
model cannot price them reliably.
\widetablebox{figures/s3_table_products.tex}{Table 5 --- Building blocks,
products and roles. The last column is the true rate sensitivity.}

\subheading{How the allocation was chosen}
\begin{enumerate}[leftmargin=1.5em, itemsep=2pt]
  \item Every allocation in 5\% steps was checked against the client's rules:
  \textbf{40,840} pass.
  \item A quick screen on expected return and risk kept the best 40, and each
  was simulated over 5,000 years.
  \item \textbf{Option A} follows a rule fixed in advance: the highest median
  return without a higher chance of breaching the drawdown limit than today.
  \textbf{Option B} follows a safety-first rule added after seeing A's
  result: the lowest risk for at least today's return.
  \item Both were validated on \textbf{5,000 new simulated years} the choice
  never saw, to remove the advantage of picking the best of many.
\end{enumerate}

\figurebox[0.95\linewidth]{figures/s3_fig_frontier.pdf}{Figure 4 --- Every
allowed allocation by expected risk and return. Today's allocation already
sits near the upper edge; options A and B are close to it.}

\tablebox{figures/s3_table_validation.tex}{Table 6 --- Validation on 5,000 new
simulated years.}

\tablebox{figures/s3_table_decision.tex}{Table 7 --- Choosing between the two
options.}

\subheading{Why option B}
On the new years, option B beat option A by 2.6 standard errors, needs only
10 points of trading instead of 25, keeps the 15\% cash buffer the risk
triggers rely on, and has the smallest rate risk. Against today, its return
gain is small (+0.11 point a year) but statistically real, because every
allocation was tested on the same simulated years.

\widetablebox{figures/s3_table_allocation.tex}{Table 8 --- Recommended
allocation (option B), products, trades and risk shares.}

\figurebox[0.85\linewidth]{figures/s3_fig_allocation.pdf}{Figure 5 --- Today
vs.\ recommended allocation.}

\subheading{What it does and does not fix}
\begin{itemize}[leftmargin=1.5em, itemsep=2pt]
  \item \textbf{Fixes:} rate risk, by replacing the most rate-sensitive sector
  (Real Estate) with two that gain when rates rise (Financials and Energy).
  \item \textbf{Does not fix:} Technology still carries about half the risk.
  Option A reduces it to 43\%, but only by putting 40\% in Healthcare and
  running cash at the 5\% minimum. The risk triggers remain the stronger tool
  for the drawdown limit.
\end{itemize}

% ================= 4. Reconciliation =================
\heading{4. Faster, Reliable Trade Reporting}
Today an analyst checks each trade by hand against the custodian's records.
I built and tested a rule-based reconciliation engine on the year's 3,000
trades. The custodian file contained 270 real breaks (wrong price, quantity
or date, missing or duplicate bookings) and 30 decoys: the same trade ID
written differently, which is not a break. The engines never saw the answer
key; it was used only to score them.

\tablebox{figures/s4_table_scores.tex}{Table 9 --- A naive spreadsheet lookup
vs.\ the rule-based engine.}

\figurebox[0.9\linewidth]{figures/s4_fig_scores.pdf}{Figure 6 --- Real breaks
found, missed, and false alarms.}

\subheading{Why the naive lookup fails, and what the engine does instead}
\tablebox{figures/s4_table_naive_errors.tex}{Table 10 --- Errors of the naive
lookup by cause.}
The engine applies four rules: standardise trade IDs, count bookings per
trade, allow prices to differ by up to 10 basis points for rounding, and give
each trade one break type in order of severity. Figure 7 shows the tolerance
is not fragile: there are no errors anywhere between 2 and 75 bp.

\figurebox[0.88\linewidth]{figures/s4_fig_tolerance.pdf}{Figure 7 --- Precision
and recall for different price tolerances.}

\subheading{The exception report and the time saved}
Each break comes with a priority (by money at stake, with missing and
duplicate bookings always high) and a suggested action.
\tablebox{figures/s4_table_exceptions.tex}{Table 11 --- The top of the
exception report (full list in \texttt{outputs/exceptions\_report.csv}).}
\tablebox{figures/s4_table_workload.tex}{Table 12 --- Analyst workload,
assuming 20 minutes per trade checked by hand and 10 minutes per flagged
trade.}

A score of 100\% is not the headline: the planted breaks are cleanly
separated from the rounding noise, whereas real reconciliations also involve
partial fills, corporate actions, currencies and split accounts. The
meaningful results are the comparison with the naive lookup, the safe range
for the tolerance, and the 96\% reduction in analyst time.

% ================= Next steps =================
\heading{Next Steps}
\begin{enumerate}[leftmargin=1.5em, itemsep=2pt]
  \item Agree with the client whether the 10\% drawdown limit is a hard
  limit; this decides whether the risk triggers are adopted.
  \item Implement the allocation change (Financials +5, Energy +5, Real
  Estate \textminus10) at the next monthly rebalance.
  \item Pilot the reconciliation engine alongside the manual check for one
  month, then switch the manual check to exceptions only.
  \item Review the rate scenarios quarterly instead of relying on estimated
  sensitivities.
\end{enumerate}

% ================= Appendix =================
\newpage
\heading{Appendix --- How This Was Built}
\begin{itemize}[leftmargin=1.5em, itemsep=3pt]
  \item \textbf{Synthetic data} (\texttt{src/generate\_data.py}, fixed
  seed): weekly sector returns from a factor model (market move, rate change,
  event shocks and noise for each sector), a mean-reverting VIX, a policy rate
  with scripted hikes and cuts, and 3,000 trades recorded by two systems with
  planted breaks and a ground-truth file.
  \item \textbf{Simulation} (\texttt{src/portfolio\_simulation.py}, shared by
  Stages 2 and 3): 5,000 random years from the same model, with weekly
  chances of a VIX jump (8\%), a $\pm$25 bp rate move (6\%) and a $\pm$4\%
  sector shock (3\%). Management approaches are applied week by week using
  only the previous week's information; trading costs are 5 bp.
  \item \textbf{Allocation search} (Stage 3): all 5\%-step allocations
  within the client's rules; screen on expected return and volatility;
  full simulation of the shortlist; validation on a fresh set of simulated
  years; paired comparison (same years for every allocation) to measure
  whether differences exceed simulation noise.
  \item \textbf{Reconciliation:} a naive exact lookup and a rule-based engine
  (file \texttt{src/reconciliation.py}), scored against the ground truth.
  \item \textbf{Limits:} the market model is the author's; real markets have
  fatter tails, changing correlations and regime shifts. Results compare
  approaches under the same assumptions and are not forecasts.
  \item \textbf{Reproduce:} regenerate the data (optional, it is committed),
  run the four notebooks in order, then compile this report:\par
  {\small\texttt{python src/generate\_data.py}\par
  \texttt{jupyter nbconvert -{}-execute -{}-inplace 0*.ipynb}\par
  \texttt{pdflatex Capstone\_Report.tex}}
\end{itemize}
\end{document}
